% ============================================================
%  Capítulo 2 — Fundamentação Teórica
%  Introdução do capítulo (texto antes da primeira \section)
% ============================================================

\chapter{Fundamentação Teórica}
\label{Cap:Fund_Teo}

Este capítulo apresenta o referencial teórico que sustenta o presente trabalho. A construção do arcabouço conceitual parte da teoria de mercados financeiros --- em especial da Hipótese de Mercados Adaptativos \cite{lo2004}, que reconhece o papel das emoções e dos comportamentos coletivos na formação dos preços --- e avança pelas bases técnicas que viabilizam a análise proposta: o fenômeno do \textit{Big Data} e a emergência das mídias sociais como fonte de dados financeiros \cite{laney2001,bollen2011}, os fundamentos do Processamento de Linguagem Natural \cite{caseli2024cap1,jm3} e as abordagens metodológicas da Análise de Sentimento aplicadas ao domínio financeiro em língua portuguesa
\cite{kearney2014,santos2023finbert}. O capítulo encerra com uma revisão dos principais trabalhos relacionados, situando esta pesquisa na literatura nacional e internacional da área.

% ============================================================
%  Seção 2.1 — Mercado Financeiro Brasileiro
%  Capítulo 2 — Fundamentação Teórica
% ============================================================

\section{Mercado Financeiro Brasileiro}
\label{sec:mercado_financeiro}

O mercado financeiro brasileiro é organizado em torno da B3 (Brasil, Bolsa, Balcão), que concentra a negociação de ações, derivativos, títulos de renda fixa e câmbio em uma única infraestrutura integrada. Em termos de capitalização e liquidez, a B3 figura entre as principais bolsas de valores da América Latina, sendo amplamente utilizada como referência para estudos empíricos sobre o comportamento de ativos em mercados emergentes \cite{costajr1994,yoshinaga2012}. A dinâmica desse mercado é influenciada por um conjunto heterogêneo de fatores --- fundamentos macroeconômicos, decisões de política monetária, eventos políticos e, crescentemente, o fluxo contínuo de informações e opiniões veiculadas por meios digitais --- o que torna a análise de sentimento uma ferramenta potencialmente relevante para compreender os movimentos coletivos dos investidores.

% ============================================================
%  Subseção 2.1.1 — Hipótese de Mercados Eficientes e
%                   Hipótese de Mercados Adaptativos
% ============================================================

\subsection{Hipótese de Mercados Eficientes e Hipótese de Mercados Adaptativos}
\label{subsec:hme_hma}

A Hipótese de Mercados Eficientes (HME), formulada por \citeauthor{fama1970} (\citeyear{fama1970}) e consolidada em sua versão revisada por \citeauthor{fama1991} (\citeyear{fama1991}), postula que os preços dos ativos financeiros refletem, a qualquer momento, toda a informação disponível relevante. Em sua forma semiforte, a HME implica que nenhuma análise de informações públicas --- incluindo notícias, relatórios ou declarações de agentes econômicos --- seria capaz de gerar retornos sistematicamente superiores ao mercado, uma vez que essas informações já estariam incorporadas aos preços de forma quase instantânea \cite{fama1970,grossman1980}.

A HME, contudo, foi progressivamente desafiada por evidências empíricas de anomalias de mercado --- como sobrereação e subreação a eventos, efeito momentum e reversão à média --- que não se encaixam no paradigma de plena eficiência informacional \cite{costajr1994}. No contexto brasileiro, \citeauthor{costajr1994} (\citeyear{costajr1994}) documentaram evidências de sobrereação no mercado acionário nacional, sugerindo que investidores locais respondem de forma excessiva a notícias recentes, padrão consistente com a influência de fatores comportamentais e emocionais nas decisões de investimento.

Em resposta a essas limitações, \citeauthor{lo2004} (\citeyear{lo2004,lo2005}) propôs a Hipótese de Mercados Adaptativos (HMA), que reinterpreta a eficiência de mercado sob uma perspectiva evolucionária. Segundo a HMA, os agentes financeiros não são racionais no sentido neoclássico estrito, mas adaptam seu comportamento continuamente em resposta ao ambiente de mercado, incorporando heurísticas, emoções e aprendizado social em suas decisões \cite{lo2004}. Essa perspectiva é particularmente relevante para o presente trabalho: se os preços dos ativos refletem não apenas informações objetivas, mas também o estado emocional e as expectativas coletivas dos agentes --- conforme previsto pela HMA --- então o sentimento expresso em publicações do X/Twitter constitui um sinal informacional potencialmente relevante para compreender o humor do mercado financeiro brasileiro \cite{lo2004,bollen2011,santos2023finbert}. Não por acaso, \citeauthor{santos2023finbert} (\citeyear{santos2023finbert}) invocam explicitamente a HMA como fundamento teórico do FinBERT-PT-BR, situando a análise de sentimento como ferramenta de captura das ``emoções refletidas nos textos'' que influenciam as decisões dos agentes no mercado adaptativo.

\input{USPSC-TA-Textual/USPSC-Cap2-Fundamentacao_Teorica/USPSC-Cap2-Fundamentacao_Teorica_2.2-Big_Data_e_Midias_Sociais_em_Financas.tex}

\input{USPSC-TA-Textual/USPSC-Cap2-Fundamentacao_Teorica/USPSC-Cap2-Fundamentacao_Teorica_2.3-Processamento_de_Linguagem_Natural.tex}

\input{USPSC-TA-Textual/USPSC-Cap2-Fundamentacao_Teorica/USPSC-Cap2-Fundamentacao_Teorica_2.4-Analise_de_Sentimento.tex}

\input{USPSC-TA-Textual/USPSC-Cap2-Fundamentacao_Teorica/USPSC-Cap2-Fundamentacao_Teorica_2.5-Trabalhos_Relacionados.tex}
